\documentclass[11pt]{article}
\usepackage[margin=1in]{geometry}
\usepackage{amsmath,amssymb,amsthm,mathtools}
\usepackage[T1]{fontenc}
\usepackage{hyperref}
\hypersetup{colorlinks=true,linkcolor=blue,urlcolor=blue,citecolor=blue}
\newtheorem{theorem}{Theorem}
\newtheorem{lemma}[theorem]{Lemma}
\newcommand{\lcmop}{\operatorname{lcm}}
\newcommand{\Sph}{\widehat R_G}
\title{A determinant criterion for extending generalized splines\\
\large MathDB 372959: a reduction to ordinary graph splines}
\author{Principia Math In Progress}
\date{18 September 2026}
\begin{document}
\maketitle

\paragraph{Task 1: the original problem, solved.}
Conjecture 1.3 of Dilaver--Alt\i nok~\cite{DA} asks whether every basis
of the extending spline module has determinant associated to its
specified element $\widehat Q_G$. We prove this for every GCD domain,
without the coprimality or PID restrictions in the source. The input
is the existing ordinary-spline theorem of Fi\c sekci--Sar\i o\u glan~\cite{FS};
the reduction and the exact identification of its determinant factors
are given below.

Let $G$ have ordered vertices $v_1,\ldots,v_n$. In the source's setting,
vertex and edge modules are $m_iR$ and $R/r_eR$, with the natural quotient
maps. Thus
\[
 \Sph=\{f\in R^n:m_i\mid f_i\text{ for every }i,
       \quad r_{ij}\mid f_i-f_j\text{ for every }v_iv_j\in E(G)\}.
\]
The labels are nonzero. In fact, the existence of an $n$-element basis
already excludes a zero vertex label or a zero label on an edge between
distinct vertices: either imposes a nontrivial linear relation over the
fraction field, giving rank less than $n$. Loops impose no condition and
may be omitted. Empty graphs and disconnected graphs are allowed.
All gcds, lcms and equalities between their products are understood up
to units; the lcm of an empty family is $1$.

For distinct vertices put
\[
 d_{ij}=\lcmop_{P:v_i\leadsto v_j\text{ in }G}\gcd(r_e:e\in P).
\]
One may use all trails, as in the source, or just simple paths: removing
a closed subwalk increases the path gcd in the divisibility order, and
every simple path is a trail. If there is no path, set $d_{ij}=1$.
Definition 3.2 of~\cite{DA} is precisely
\begin{equation}\label{eq:Q}
 q_i=\lcmop\bigl(m_i,\{\gcd(m_j,d_{ij}):j>i\},
                         \{d_{is}:s<i\}\bigr),\qquad
 \widehat Q_G=\prod_{i=1}^n q_i.
\end{equation}

\begin{theorem}\label{thm:main}
For any $F_1,\ldots,F_n\in\Sph$ over a GCD domain $R$, these splines
form an $R$-basis if and only if
\[
 \det(F_1,\ldots,F_n)=u\widehat Q_G\qquad(u\in R^\times).
\]
In particular, the conjectured necessary condition holds in full.
\end{theorem}

Add a new vertex $v_0$, join it to each $v_i$ by an edge labelled $m_i$,
and retain every edge and label of $G$. Denote the resulting ordinary
edge-labelled graph by $H$, ordered $v_0<v_1<\cdots<v_n$.
Let $S_H\subseteq R^{n+1}$ be its ordinary spline module.

\begin{lemma}\label{lem:split}
The map
\[
 R\oplus\Sph\longrightarrow S_H,\qquad
 (a,f)\longmapsto(a,a+f_1,\ldots,a+f_n)
\]
is an isomorphism. Moreover, the columns
$\mathbf1,(0,F_1),\ldots,(0,F_n)$ have determinant
$\det(F_1,\ldots,F_n)$ in the ascending row order.
\end{lemma}
\begin{proof}
For $g\in S_H$, the spoke conditions give $m_i\mid g_i-g_0$,
and the old edge conditions are unchanged by subtracting $g_0$ from
each coordinate. The inverse is therefore
$g\mapsto(g_0,(g_i-g_0)_{i=1}^n)$.
The asserted matrix has first row $(1,0,\ldots,0)$; expand along that row.
\end{proof}

\begin{lemma}\label{lem:factors}
Let $L_0=1$, and for $i\geq1$ let $L_i$ be the lcm of path gcds in $H$
from $v_i$ to any earlier vertex in $\{v_0,\ldots,v_{i-1}\}$.
Then $L_i=q_i$ up to a unit, and the ordinary determinant element
$Q_H=\prod_{i=0}^n L_i$ equals $\widehat Q_G$ up to a unit.
\end{lemma}
\begin{proof}
The direct spoke gives $m_i\mid L_i$. Every old path to $v_s$, $s<i$,
is still a path in $H$, giving $d_{is}\mid L_i$.
For $j>i$, append the spoke $v_jv_0$ to each simple path in $G$
from $v_i$ to $v_j$. These are paths to an earlier vertex. Gcd/lcm
distributivity in a GCD domain gives
\[
 \gcd(m_j,d_{ij})
 =\lcmop_{P:v_i\leadsto v_j\text{ in }G}
       \gcd(m_j,\{r_e:e\in P\})\mid L_i.
\]
The same assertion is vacuous when the path family is empty.
Consequently $q_i\mid L_i$.

Conversely, take any simple path in $H$ from $v_i$ to an earlier vertex
and truncate it at its first visit to $\{v_0,\ldots,v_{i-1}\}$.
The gcd of the original path divides that of this prefix.
If the prefix ends at $v_s$ with $1\leq s<i$, it stays entirely in $G$,
so its gcd divides $d_{is}$. If it ends at $v_0$, its last edge is a
spoke $v_jv_0$. A direct spoke has gcd $m_i$. Otherwise $j>i$, and
the preceding segment lies in $G$, so the prefix gcd divides
$\gcd(m_j,d_{ij})$. Each case gives a divisor of $q_i$.
Taking the lcm over all paths proves $L_i\mid q_i$.
\end{proof}

\begin{proof}[Proof of Theorem~\ref{thm:main}]
The ordinary determinant criterion~\cite[Theorem 4.5, p.~367]{FS}
states that on any finite edge-labelled graph over a GCD domain,
a set of as many splines as vertices is a basis exactly when its
determinant is a unit times the product of these earlier-vertex path
factors. This is Theorem 5.3 in the arXiv version of that paper.

By Lemma~\ref{lem:split}, the proposed $F_1,\ldots,F_n$ form a basis of
$\Sph$ exactly when $\mathbf1,(0,F_1),\ldots,(0,F_n)$ form a basis of
$S_H$. Apply the cited theorem on $H$, use its block determinant,
and then use Lemma~\ref{lem:factors}. Reversing the row order, as in
the primary source, changes only a sign, which is a unit.
\end{proof}

\paragraph{Scope and verification.}
No flow-up basis, Bezout identity, prime factorization, or Noetherian
assumption is used. In particular, gcd $1$ is never mistaken for
comaximality. The source's PID and pairwise-coprime cases are existing
results, and the ordinary-spline theorem is explicitly credited.
The accompanying exact verifier compares independently enumerated
path factors, counts finite residue lattices over $\mathbb Z$, and
checks polynomial-labelled examples. These finite computations are
regression checks; the two lemmas prove the assertion for all graphs
and GCD domains.

\paragraph{Task 2 (proposed, not solved here).}
Find an order-independent expression for $\widehat Q_G$ in terms of
rooted spanning forests of $G$ that remains valid when $\Sph$ is not free,
and compare its arithmetic complexity with direct path enumeration.

\paragraph{Task 3 (proposed, not solved here).}
For a cycle $G$ over $k[x,y,z]$, with every vertex and edge label a
nonzero linear form, characterize exactly when $\Sph$ is free in terms
of the linear dependencies among the labels.

\paragraph{Task 4 (proposed, not solved here).}
Over a Dedekind domain, replace the principal vertex modules $m_iR$
by invertible ideals $I_i$. Determine the determinant ideal and the
Steinitz class of the resulting spline module from the labelled graph.

\paragraph{Task 5 (proposed, not solved here).}
Replace the edge congruence by
$a_{ij}f_i-b_{ij}f_j\in r_{ij}R$, where $a_{ij},b_{ij}$ are units.
Characterize when multiplication of vertex coordinates by units turns
this system into an ordinary augmented graph, and obtain the determinant
criterion in the remaining cases with nontrivial cycle holonomy.

\begin{thebibliography}{9}
\bibitem{DA} G. Dilaver and S. Alt\i nok,
\emph{Basis Criteria for Extending Generalized Splines},
\href{https://arxiv.org/abs/2602.04440v1}{arXiv:2602.04440v1} (2026),
Conjecture 1.3 and Definition 3.2.
\bibitem{FS} S. Fi\c sekci and S. Sar\i o\u glan,
\emph{Basis condition for generalized spline modules},
J. Algebraic Combin. \textbf{59} (2024), 359--369,
\href{https://doi.org/10.1007/s10801-023-01290-y}{doi:10.1007/s10801-023-01290-y};
\href{https://arxiv.org/abs/2301.12529}{arXiv:2301.12529}.
\end{thebibliography}
\end{document}
